% Conclusion

\section{Conclusion}
\label{sec:conclusion}

We presented OmniFlow, an adaptive replay system for mobile GUI agents that
turns past successful executions into reusable \texttt{Functions}.
The core design separates reusable execution knowledge from live execution
control.   \textsc{Flow} applies those Functions
on the live device by observing the current state, checking pending actions,
transferring recorded targets to current screens, and falling back when the
evidence is insufficient.

This representation/execution split changes the cost model of repeated GUI serving.  A
successful run can write back a reusable \texttt{Function}; later compatible
users, devices, or task variants can execute the matched Function without
re-planning it.  Cache hits skip online reasoning for the matched Function,
while misses and unsafe transfers return compact page and execution context to
the host agent.  The online cost is therefore amortized over newly discovered
reusable Functions rather than over every user, goal, and repetition.

OmniFlow currently focuses on Android GUI execution and still depends on
sufficient state similarity for safe reuse.  Large UI redesigns, highly dynamic
content regions, and true cold-start goals require fallback to the host agent
and later write-back.  Its adaptation also has a clear scope: Flow can transfer
source actions across current GUI states when structure, geometry, and effect
evidence support a confident target, while open semantic decisions can remain
VLM, script, or host-agent boundaries.  This makes OmniFlow a practical
adaptive replay layer: executable, privacy-safe \texttt{Function}s are shared,
transferred, checked, and refined from future run logs.
